\subsection{正则叶片与适配度规}

可积分布给出局部叶片以后，度规还须确定法向距离、切向重标记与积分测度。

\begin{definition}[正则叶状分解]
\label{def:ghs-regular-foliation}
设 $U\subset M$ 是开区域，$I\subset\RR$ 是区间，$\Sigma$ 是 $(d-1)$ 维流形。若存在微分同胚 $F:I\times\Sigma\longrightarrow U$，则 $\Sigma_t=F(\{t\}\times\Sigma)$ 称为一族正则叶。函数 $t:U\to I$ 在每点满足 $\dd t\neq 0$，因此每个 $\Sigma_t$ 都是正则超曲面。
\end{definition}

此定义的用处是将“量子化面”与特定的球面分开。譬如平直空间去掉原点后有 $U\simeq\RR\times S^{d-1}$；而一般紧致流形可以没有全局非临界实函数，因此不存在这样的全局展开。临界点、拓扑变化与边界必须另行处理，不能一律塞进所谓的径向时间。

\begin{proposition}[叶状度规的分解]
\label{prop:ghs-foliation-metric}
在 $I\times\Sigma$ 上任取适配坐标 $(t,y^a)$，正定度规唯一写成
\begin{equation}\label{eq:ghs-foliation-adm}
 g=N^2\dd t^2+h_{ab}(\dd y^a+N^a\dd t)(\dd y^b+N^b\dd t),
\end{equation}
其中 $N>0$ 是 法向进动系数，$N^a$ 是 切向位移，$h_{ab}$ 是 $\Sigma_t$ 上的诱导度规。体积测度为
\begin{equation}\label{eq:ghs-foliation-volume}
 \dd\mu_g=N\sqrt{h}\,\dd t\,\dd^{d-1}y.
\end{equation}
\end{proposition}
\begin{proof}
按 $t$ 和 $y$ 分块写度规矩阵。由于 $h_{ab}$ 正定可逆，令 $N_a=g_{ta}$、$N^a=h^{ab}N_b$，再令 $N^2=g_{tt}-h^{ab}N_aN_b$。正定性和 Schur 补定理给出 $N^2>0$。展开 \cref{eq:ghs-foliation-adm} 得到原矩阵全部分量。分块行列式公式
\[
 \det g=\det h\,(g_{tt}-N_ah^{ab}N_b)=N^2\det h
\]
便证明测度公式。唯一性来自矩阵分块分解的唯一性。
\end{proof}

把 $t$ 增加 $\dd t$ 时，坐标向量 $\partial_t$ 一般不是单位法向。单位法向和法向微分为
\[
 n=N^{-1}(\partial_t-N^a\partial_a),\qquad
 \nabla_n\varphi=N^{-1}(\partial_t-N^a\partial_a)\varphi.
\]
这决定了标量场在两片叶之间真正改变了多少。不能在存在 切向位移 时仅用 $\partial_t\varphi$ 作为正交时间导数。
